\chapter{Glossary}
\label{ch:glossary}

\section{Purpose}

This glossary defines terms as they are used \emph{in this document} and in the
Expo Snips repository. Where a term is specific to the project it is marked
\textbf{Project}; where it is an ecosystem term it is marked \textbf{General}.
Definitions are grounded in the repository code.

\section{Domain terms}

\begin{longtable}{@{}L{3.8cm}L{1.6cm}L{8.8cm}@{}}
\caption{Glossary of domain terms}\\
\toprule
\textbf{Term} & \textbf{Type} & \textbf{Definition} \\
\midrule
\endfirsthead
\toprule
\textbf{Term} & \textbf{Type} & \textbf{Definition} \\
\midrule
\endhead
Recipe & Project & A JSON document describing one screen: its identity, version,
supported options, dependencies, output files, optional variants, tags and
primitives. Lives in \texttt{packages/registry/data/}. \\
Template & Project & A Handlebars (\texttt{.hbs}) file that renders to a source
file in the target project. Lives in \texttt{packages/registry/templates/}. \\
Registry & Project & The published catalog of recipes: \texttt{index.json} plus
one JSON document per recipe. Static, hosted on GitHub Pages. \\
Index entry & Project & The summary of a recipe used for browsing and quick
picks: id, name, description, version, typescript flag, tags and the supported
options. \\
Catalog & Project & Collectively, all recipes. Shipped inside each client as an
offline bundle. \\
Offline bundle & Project & The inlined catalog (\texttt{bundle.json} or the
generated \texttt{registry-data.ts}) that lets a client work with no network. \\
Variant & Project & A named alternative style for a screen (for example
\texttt{glass}). Declared at the recipe level; a file maps a variant to its
template. \\
Primitive & Project & A small reusable UI recipe (for example \texttt{ui/button},
\texttt{ui/input}, \texttt{ui/utils}) required when \texttt{ui=shadcn}, resolved
recursively and skipped if already present. \\
Generation & Project & The act of turning a recipe plus options into written
files. \\
Generation options & Project & The resolved choices that shape output:
\texttt{styling}, \texttt{router}, \texttt{ui}, \texttt{icons} and
\texttt{typescript}. \\
Detection & Project & Inferring options from the target project by probing
\texttt{package.json} and \texttt{tsconfig.json}. \\
Manifest & Project & \texttt{.expo-snips.json} in the target project, recording
which recipes (and variants) have been generated, their versions, options and
files. \\
Upsert key & Project & The identity used to record a generated screen: recipe id
plus variant. \\
\bottomrule
\end{longtable}

\section{Ecosystem and technology terms}

\begin{longtable}{@{}L{3.8cm}L{1.6cm}L{8.8cm}@{}}
\caption{Glossary of technology terms}\\
\toprule
\textbf{Term} & \textbf{Type} & \textbf{Definition} \\
\midrule
\endfirsthead
\toprule
\textbf{Term} & \textbf{Type} & \textbf{Definition} \\
\midrule
\endhead
Expo & General & A framework and toolchain for building React Native applications;
the runtime target of every generated screen. \\
React Native & General & A framework for building native mobile apps with React;
the platform target. \\
Expo Router & General & File-based routing for Expo apps; the default router
option. \\
React Navigation & General & An alternative navigation library; the other router
option. \\
NativeWind & General & A Tailwind-style utility-class system for React Native; the
default styling option. \\
Tamagui & General & A UI and styling library; listed as a styling option but not
yet implemented in any recipe. \\
shadcn & General/Project & Here, a set of copy-in UI primitives (Button, Input,
and a \texttt{cn} utility) and design tokens; selectable as a UI option and
backed by primitive recipes. \\
lucide / lucide-react-native & General & The default icon library. \\
Vector Icons & General & An icon library listed as an option but not implemented in
any recipe. \\
safe area & General & The screen region not obscured by device notches, status
bars or home indicators; generated screens handle it via
\texttt{react-native-safe-area-context}. \\
Handlebars & General & The template engine used by the shared core. \\
esbuild & General & The bundler used for the VS Code extension. \\
vsce / VSIX & General & The VS Code packaging tool and the packaged extension
format. \\
pnpm workspace & General & The monorepo tooling; packages are linked from
\texttt{packages/*}. \\
Monorepo & General & A single repository containing multiple packages. \\
Longest-common-subsequence diff & General & The algorithm behind the update
preview; compares two files line by line. \\
Golden-file test & General & A test that compares generated output against a
stored expected file; recommended for this project. \\
\bottomrule
\end{longtable}

\section{Documentation and status terms}

\begin{longtable}{@{}L{3.8cm}L{1.6cm}L{8.8cm}@{}}
\caption{Glossary of documentation terms}\\
\toprule
\textbf{Term} & \textbf{Type} & \textbf{Definition} \\
\midrule
\endfirsthead
\toprule
\textbf{Term} & \textbf{Type} & \textbf{Definition} \\
\midrule
\endhead
IMPLEMENTED & This doc & The behaviour is present in the code, with evidence. \\
INFERRED & This doc & A reasonable conclusion drawn from code or structure, not
stated explicitly. \\
RECOMMENDED & This doc & A proposed improvement; not present in the code. \\
UNKNOWN & This doc & Cannot be determined from the repository. \\
NOT APPLICABLE & This doc & The concept does not exist in this product. \\
OPEN / ACCEPTED / PASS & This doc & Status of a finding: unaddressed, knowingly
accepted, or satisfied. \\
Finding / Severity / Evidence / Impact / Recommendation & This doc & The fixed
format used for security findings in Chapter~\ref{ch:security}. \\
BRD & This doc & Business Requirements Document. \\
PRD & This doc & Product Requirements Document. \\
SEO & This doc & Search Engine Optimisation. \\
AEO & This doc & Answer Engine Optimisation. \\
GEO & This doc & Generative Engine Optimisation. \\
\bottomrule
\end{longtable}

\section{Abbreviations used in this document}

\begin{longtable}{@{}L{2.6cm}L{11.6cm}@{}}
\caption{Abbreviations}\\
\toprule
\textbf{Abbreviation} & \textbf{Expansion} \\
\midrule
\endfirsthead
\toprule
\textbf{Abbreviation} & \textbf{Expansion} \\
\midrule
\endhead
CLI & Command-Line Interface. \\
CI / CD & Continuous Integration / Continuous Delivery. \\
CORS & Cross-Origin Resource Sharing. \\
DX & Developer Experience. \\
IaC & Infrastructure as Code. \\
JSON-LD & JavaScript Object Notation for Linked Data. \\
LCS & Longest Common Subsequence. \\
N/A & Not Applicable. \\
PASS & A security control that is satisfied. \\
RBAC & Role-Based Access Control (not applicable to this project). \\
SPDX & Software Package Data Exchange (a licence identifier standard). \\
VSIX & VS Code extension package format. \\
\bottomrule
\end{longtable}
